% Options for packages loaded elsewhere
\PassOptionsToPackage{unicode}{hyperref}
\PassOptionsToPackage{hyphens}{url}
\documentclass[
  12pt,
]{ctexart}
\usepackage{xcolor}
\usepackage{amsmath,amssymb}
\setcounter{secnumdepth}{-\maxdimen} % remove section numbering
\usepackage{iftex}
\ifPDFTeX
  \usepackage[T1]{fontenc}
  \usepackage[utf8]{inputenc}
  \usepackage{textcomp} % provide euro and other symbols
\else % if luatex or xetex
  \usepackage{unicode-math} % this also loads fontspec
  \defaultfontfeatures{Scale=MatchLowercase}
  \defaultfontfeatures[\rmfamily]{Ligatures=TeX,Scale=1}
\fi
\usepackage{lmodern}
\ifPDFTeX\else
  % xetex/luatex font selection
\fi
% Use upquote if available, for straight quotes in verbatim environments
\IfFileExists{upquote.sty}{\usepackage{upquote}}{}
\IfFileExists{microtype.sty}{% use microtype if available
  \usepackage[]{microtype}
  \UseMicrotypeSet[protrusion]{basicmath} % disable protrusion for tt fonts
}{}
\makeatletter
\@ifundefined{KOMAClassName}{% if non-KOMA class
  \IfFileExists{parskip.sty}{%
    \usepackage{parskip}
  }{% else
    \setlength{\parindent}{0pt}
    \setlength{\parskip}{6pt plus 2pt minus 1pt}}
}{% if KOMA class
  \KOMAoptions{parskip=half}}
\makeatother
\usepackage{longtable,booktabs,array}
\usepackage{calc} % for calculating minipage widths
% Correct order of tables after \paragraph or \subparagraph
\usepackage{etoolbox}
\makeatletter
\patchcmd\longtable{\par}{\if@noskipsec\mbox{}\fi\par}{}{}
\makeatother
% Allow footnotes in longtable head/foot
\IfFileExists{footnotehyper.sty}{\usepackage{footnotehyper}}{\usepackage{footnote}}
\makesavenoteenv{longtable}
\setlength{\emergencystretch}{3em} % prevent overfull lines
\providecommand{\tightlist}{%
  \setlength{\itemsep}{0pt}\setlength{\parskip}{0pt}}
\usepackage{bookmark}
\IfFileExists{xurl.sty}{\usepackage{xurl}}{} % add URL line breaks if available
\urlstyle{same}
\hypersetup{
  hidelinks,
  pdfcreator={LaTeX via pandoc}}

\author{SCX}
\date{2026}

\begin{document}

\section{SCX
核心数学定义速查}\label{scx-ux6838ux5fc3ux6570ux5b66ux5b9aux4e49ux901fux67e5}

\begin{quote}
内容来源：AI 对话（2026-06-25），待形式化整理到 \texttt{theory/} 目录
\end{quote}

\begin{center}\rule{0.5\linewidth}{0.5pt}\end{center}

\subsection{1. 问题设定}\label{ux95eeux9898ux8bbeux5b9a}

设： - 输入空间 X，标签空间 Y - 未知真实标注函数 f*: X → Y - M 个专家 E
= \{f\_1, f\_2, \ldots, f\_M\} - 状态映射 s: X → S（或软分配
p(s\textbar x)） - 表示映射 φ: X → R\^{}d

\subsection{2. 核心对象}\label{ux6838ux5fc3ux5bf9ux8c61}

\subsubsection{Definition 1: State-Conditioned Expert
Risk}\label{definition-1-state-conditioned-expert-risk}

\begin{verbatim}
R_m(s) = E_{x~P(·|s)} [ ℓ(f_m(x), f*(x)) ]
\end{verbatim}

专家的风险不是全局量 R\_m，而是状态条件量 R\_m(s)。

\subsubsection{Definition 2: SCX
Reliability}\label{definition-2-scx-reliability}

\begin{verbatim}
SCX_m(s) = P( ℓ(f_m(x), y) < τ | x ∈ s )
\end{verbatim}

即专家 m 在状态 s 下做出可接受预测的概率。

\subsubsection{Definition 3: State Data
Value}\label{definition-3-state-data-value}

\begin{verbatim}
V(s) = r̄(s) · ρ(s) · L(s) · [1 - D(s)] · max_m SCX_m(s)
\end{verbatim}

其中： - \texttt{r̄(s)\ =\ E{[}ℓ(f(x),\ y)\ \textbar{}\ s{]}} ---
当前模型在该状态的平均误差 - \texttt{ρ(s)\ =\ P(x\ ∈\ s)} ---
状态出现概率/质量 - \texttt{L(s)} --- 可学习性（状态内一致、非噪声） -
\texttt{D(s)} --- 冗余度（已充分覆盖程度）

\subsubsection{Definition 4: Learnability
Score}\label{definition-4-learnability-score}

\begin{verbatim}
L(s) = C(s) · [1 - N(s)]
\end{verbatim}

\begin{itemize}
\tightlist
\item
  C(s): 状态内一致性（标签一致、专家预测一致）
\item
  N(s): 噪声分数
\end{itemize}

\subsubsection{Definition 5: Noise Score (per
sample)}\label{definition-5-noise-score-per-sample}

\begin{verbatim}
NoiseScore(x_i) = r_i · (1/ρ(s_i)+ε) · [1 - C(s_i)]
\end{verbatim}

\begin{itemize}
\tightlist
\item
  高误差 + 低密度 + 低一致性 → 更像噪声
\item
  高误差 + 高密度 + 高一致性 → 更像可学习困难状态
\end{itemize}

\subsection{3. 数据四分类}\label{ux6570ux636eux56dbux5206ux7c7b}

\begin{longtable}[]{@{}
  >{\raggedright\arraybackslash}p{(\linewidth - 4\tabcolsep) * \real{0.3333}}
  >{\raggedright\arraybackslash}p{(\linewidth - 4\tabcolsep) * \real{0.3333}}
  >{\raggedright\arraybackslash}p{(\linewidth - 4\tabcolsep) * \real{0.3333}}@{}}
\toprule\noalign{}
\begin{minipage}[b]{\linewidth}\raggedright
类别
\end{minipage} & \begin{minipage}[b]{\linewidth}\raggedright
条件
\end{minipage} & \begin{minipage}[b]{\linewidth}\raggedright
动作
\end{minipage} \\
\midrule\noalign{}
\endhead
\bottomrule\noalign{}
\endlastfoot
\textbf{Valuable} & r̄(s)↑, ρ(s)↑, C(s)↑, D(s)↓ & acquire \\
\textbf{Redundant} & r̄(s)↓, ρ(s)↑, Coverage(s)↑ & skip \\
\textbf{Noisy} & r(x)↑, ρ(s)↓, C(s)↓ & downweight/discard \\
\textbf{Expert-dependent} & ∃m: R\_m(s) ≪ R\_n(s) for n≠m & route to
expert m \\
\end{longtable}

\subsection{4. 专家路由}\label{ux4e13ux5bb6ux8defux7531}

\begin{verbatim}
m*(x) = arg min_m Σ_s γ_s(x) · R_m(s) + λ · C_m
\end{verbatim}

或权重形式：

\begin{verbatim}
w_m(x) ∝ exp( -α · Σ_s γ_s(x) · R̂_m(s) )
\end{verbatim}

\subsection{5. State-Wise Acquisition}\label{state-wise-acquisition}

\begin{verbatim}
s* = arg max_s  r̄(s) · ρ(s) · L(s) · [1 - D(s)]
\end{verbatim}

先选高价值状态，再在状态内选代表点。

\subsection{6. 三个核心命题}\label{ux4e09ux4e2aux6838ux5fc3ux547dux9898}

\subsubsection{Proposition 1: Global Expert Ranking
Insufficiency}\label{proposition-1-global-expert-ranking-insufficiency}

若存在 s₁, s₂ 使得：

\begin{verbatim}
R_a(s₁) < R_b(s₁)  但  R_a(s₂) > R_b(s₂)
\end{verbatim}

则不存在全局最优专家排序。专家选择必须状态条件化。

\subsubsection{Proposition 2: High-Error Sampling Suboptimality under
Noise}\label{proposition-2-high-error-sampling-suboptimality-under-noise}

高误差集合 H = H\_learnable ∪ H\_noise。若：

\begin{verbatim}
P(ΔR < 0 | H_learnable) > P(ΔR < 0 | H_noise)
\end{verbatim}

则 max-residual sampling (x* = arg max r(x)) 浪费标注预算。 应改用
state-value sampling (s* = arg max r̄(s)·ρ(s)·C(s))。

\subsubsection{Proposition 3: State-Conditioned
Weighting}\label{proposition-3-state-conditioned-weighting}

最优专家权重不是全局常数，而是：

\begin{verbatim}
w_m(x) ∝ exp( -α · Σ_s γ_s(x) · R̂_m(s) )
\end{verbatim}

\subsection{7. 与材料 MLIP
的实例化}\label{ux4e0eux6750ux6599-mlip-ux7684ux5b9eux4f8bux5316}

\begin{longtable}[]{@{}ll@{}}
\toprule\noalign{}
数学对象 & 材料实例 \\
\midrule\noalign{}
\endhead
\bottomrule\noalign{}
\endlastfoot
x & 原子结构/局域环境 \\
y & DFT 能量/力 \\
f\_m & ACE/NEP/MACE expert \\
φ(x) & ACE/SOAP/描述符 \\
s & sp², sp³, 断键, 表面, Al-Ga mixed, 高应变 \\
R\_m(s) & expert 在该结构状态下的误差 \\
acquire & 补 DFT 样本 \\
relabel & DFT 仲裁 \\
downweight & 清理坏点 \\
route & 选择最合适的 expert 标注 \\
\end{longtable}

\subsection{8. 论文标题建议}\label{ux8bbaux6587ux6807ux9898ux5efaux8bae}

\begin{quote}
\textbf{SCX: State-Conditioned eXpertise for Data Valuation and
Expert-Guided Learning}
\end{quote}

中文：\textbf{SCX：用于数据价值评估与专家引导学习的状态条件专家性框架}

\end{document}
